% ============================================================
%  Housing Supply Elasticity and Transit Welfare:
%  A Replication and Extension of Severen (2021)
% ============================================================
\documentclass[12pt,a4paper]{article}

% ── Packages ────────────────────────────────────────────────
\usepackage[a4paper, margin=2.5cm]{geometry}
\usepackage{setspace}
\usepackage{microtype}
\usepackage{lmodern}
\usepackage[T1]{fontenc}
\usepackage[utf8]{inputenc}
\usepackage{mathptmx}          % Times-style body font
\usepackage{helvet}            % Helvetica for headings (used via \textsf)
\usepackage{amsmath,amssymb}
\usepackage{booktabs}
\usepackage{array}
\usepackage{tabularx}
\usepackage{multirow}
\usepackage[table]{xcolor}
\usepackage{graphicx}
\usepackage{float}
\usepackage{caption}
\usepackage[round,authoryear]{natbib}
\usepackage[hidelinks,
            colorlinks=true,
            linkcolor=NavyBlue,
            citecolor=NavyBlue,
            urlcolor=NavyBlue]{hyperref}
\usepackage{titlesec}
\usepackage{abstract}
\usepackage{enumitem}
\usepackage{parskip}

% ── Colours ─────────────────────────────────────────────────
\definecolor{NavyBlue}{HTML}{2E6DA4}
\definecolor{LightGrey}{HTML}{F2F2F2}
\definecolor{MidGrey}{HTML}{D9D9D9}
\definecolor{RuleGrey}{HTML}{AAAAAA}

% ── Section heading style ────────────────────────────────────
\titleformat{\section}
  {\large\bfseries\sffamily\color{NavyBlue}}
  {\thesection}{1em}{}[\vspace{-4pt}\textcolor{NavyBlue}{\rule{\linewidth}{0.4pt}}]

\titleformat{\subsection}
  {\normalsize\bfseries\sffamily\color{black!70}}
  {\thesubsection}{1em}{}

\titlespacing{\section}{0pt}{18pt}{6pt}
\titlespacing{\subsection}{0pt}{12pt}{4pt}

% ── Caption style ────────────────────────────────────────────
\captionsetup{
  font=small, labelfont=bf,
  justification=justified, singlelinecheck=false
}

% ── Spacing ──────────────────────────────────────────────────
\setstretch{1.15}
\setlength{\parindent}{0pt}
\setlength{\parskip}{7pt}
\setlength{\emergencystretch}{2em}

% ── Table column types ───────────────────────────────────────
\newcolumntype{L}[1]{>{\raggedright\arraybackslash}p{#1}}
\newcolumntype{C}[1]{>{\centering\arraybackslash}p{#1}}
\newcolumntype{R}[1]{>{\raggedleft\arraybackslash}p{#1}}

% ── Shaded row command ───────────────────────────────────────
\newcommand{\shaderow}{\rowcolor{LightGrey}}

% ── Footnote style for tables ────────────────────────────────
\newcommand{\tablenote}[1]{%
  \vspace{4pt}
  {\small\textit{Notes:} #1}%
}

% ============================================================
\begin{document}

% ── Title block ─────────────────────────────────────────────
\begin{titlepage}
  \vspace*{3cm}
  {\sffamily\Large\bfseries\color{NavyBlue}
    Housing Supply Elasticity and Transit Welfare:\\[6pt]
    A Replication and Extension of Severen (2021)}

  \vspace{1.2cm}
  {\large Marko Ikävalko \& Mikko Viikari}

  \vspace{0.4cm}
  {\normalsize May 2026}

  \vspace{2cm}
  \noindent\textcolor{NavyBlue}{\rule{\linewidth}{1pt}}

  \vspace{1.2cm}
  \begin{abstract}
  \noindent
  \citet{severen2021} estimates the general-equilibrium welfare effects of the
  Los Angeles Metro Rail network.  The paper finds commuting flows rose
  11--16\,\% between station-adjacent census-tract pairs and derives an
  aggregate welfare gain of roughly \$94 million per year (our replication:
  \$93.55 million), holding the city-wide inverse housing supply elasticity
  $\psi$ at its pooled estimate of 1.602 (a supply elasticity of
  $1/\psi \approx 0.62$).  We replicate the full Severen pipeline and extend the
  analysis in two directions.  First, we attempt to estimate $\psi$ separately
  for geographic sub-regions of the LA metro area (by coastal proximity, density
  quartile, and county group), motivated by the substantial spatial variation in
  local housing markets.  The Bartik shift-share instrument proves too weak for
  sub-regional identification: first-stage $F$-statistics fall well below
  conventional thresholds in every partition tried (KP $F < 5$, often $< 1$),
  against the paper's pooled KP $F \approx 12$.  Second, unable to obtain
  reliable sub-regional estimates, we vary $\psi$ over a wide range in the
  calibrated GE model.  We find the aggregate welfare estimate is essentially
  invariant to $\psi$: across a 5$\times$ range (0.75--4.0), welfare moves by
  only \$0.33 million (0.35\,\%).  This is expected for a small intervention in
  a car-dominated city, where the transit-induced demand shock near stations is
  too modest for housing supply conditions to affect aggregate welfare ---
  though, as we discuss, $\psi$ still governs the \emph{incidence} of the gains.
  ``Amplified-connectivity'' counterfactuals confirm that larger commuting
  benefits have limited impact: doubling the commuting disutility reduction
  associated with Metro connectivity adds only \$105 million per year in
  aggregate welfare.

  \medskip\noindent
  \textbf{Keywords:} transit infrastructure, general equilibrium, housing supply
  elasticity, welfare analysis, Los Angeles
  \end{abstract}

  \vfill
\end{titlepage}

\newpage
\tableofcontents
\newpage

% ============================================================
\section{Introduction}
\label{sec:intro}

Urban transit investment is among the most capital-intensive interventions
available to city governments.  Evaluating whether such investments are
welfare-improving requires a framework that accounts not only for the direct
benefits to commuters but also for the general-equilibrium adjustments in labour
and housing markets that follow.  \citet{severen2021} provides one of the most
complete such evaluations to date, applying a quantitative spatial model to the
1990--2002 build-out of the Los Angeles Metro Rail network.

The paper treats housing supply conditions as a single city-wide constant: the
inverse housing supply elasticity $\psi$, estimated at 1.602 using a Bartik
shift-share instrument (implying a housing supply elasticity of
$1/\psi = 0.624$, i.e., inelastic supply).  This is a natural baseline, but it
is a strong assumption: housing markets in the five-county LA metro area vary
enormously across coastal and inland sub-markets, across density gradients from
dense urban cores to low-density suburbs, and across counties with very
different zoning regimes.  If $\psi$ varies meaningfully across the city, and if
Metro Rail stations happen to be located in areas with systematically different
supply conditions, the aggregate welfare estimate could mask important
heterogeneity in who captures the gains from transit investment.

Our starting hypothesis was that housing supply conditions are a first-order
determinant of the incidence of transit benefits.  Near stations with inelastic
supply, demand increases are capitalised into land prices and captured by
existing landowners; near stations with elastic supply, housing quantities adjust
and the benefits are spread more broadly.  Understanding which regime prevails
matters for both efficiency and distributional assessments of transit policy.

We extend the Severen (2021) replication in two directions.  The first is an
attempt to estimate $\psi$ separately for geographic sub-regions of the LA metro
area, splitting the sample by coastal proximity, 1990 population density
quartile, and county group.  The second, motivated by the failure of the first,
is a systematic sensitivity analysis varying $\psi$ across a wide range in the
calibrated GE model.

Our key findings are as follows.  The Bartik instrument does not have enough
variation within sub-regions to separately identify $\psi$: first-stage
$F$-statistics fall well below 5 in every geographic partition we try.  The
paper's own pooled first stage is moderately strong (Kleibergen--Paap
$F = 12.3$ in the preferred specification, though between 7 and 9 in others),
so while the pooled estimate is adequately identified, there is little power to
spare for sub-sample analysis.  Turning to the sensitivity analysis, we find
that the aggregate welfare gain from Metro Rail is almost completely insensitive
to $\psi$.  Across a 5$\times$ range of $\psi$ (0.75 to 4.0), the welfare gain
moves by only \$0.33 million --- a 0.35\,\% variation around the \$93.55 million
baseline.  We show that this insensitivity is a direct consequence of the scale
of the intervention and of the closed-city welfare measure: Metro Rail is a
small shock in a car-dominated city, and the price adjustments that $\psi$
governs are largely transfers within the city rather than aggregate gains.

We also present ``amplified-connectivity'' counterfactuals showing that scaling
up the estimated commuting benefit of Metro connectivity raises welfare roughly
linearly in the multiplier.  At $2\times$ the baseline benefit, aggregate
welfare rises by \$105 million per year.  This pattern is consistent with
network coverage --- only 4\,\% of tracts lie near stations --- rather than the
magnitude of the per-pair commuting benefit being the binding constraint on
Metro Rail's welfare impact, although we note that confirming this would require
a counterfactual that expands coverage directly, which we do not simulate.

The remainder of this paper is organised as follows.  Section~\ref{sec:framework}
describes the Severen (2021) framework.  Section~\ref{sec:methods} describes our
modifications and methods.  Section~\ref{sec:results} presents results.
Section~\ref{sec:conclusion} concludes.

% ============================================================
\section{The Severen (2021) Framework}
\label{sec:framework}

\subsection{Setting and data}

\citet{severen2021} studies the commuting, labour-market, and housing-market
effects of the Los Angeles Metro Rail network, which opened in three phases: the
Blue Line (downtown LA to Long Beach) in 1990, and the Green and Red Lines in
1993--2000.  The analysis covers 2,552 census tracts in the five-county LA
metropolitan area (Los Angeles, Orange, Riverside, San Bernardino, and Ventura
counties) over the period 1990--2002.

The main data sources are: (i)~the 1990 and 2000 decennial censuses for
commuting flows between home and work tracts; (ii)~the Longitudinal
Employer-Household Dynamics (LEHD) database extended through 2015 for
robustness; (iii)~the Neighbourhood Change Database (NCDB) for pre-trend tests;
and (iv)~GIS-computed travel times between all $2{,}552\times2{,}552$ tract
pairs under observed and counterfactual transit networks.

\subsection{Commuting flow estimates}

The paper estimates a difference-in-differences specification for bilateral
commuting flows.  The treatment indicator is an indicator for whether both the
home tract and the work tract contain (or lie within 250\,m or 500\,m of) a
Metro station.  Identification rests on the parallel pre-trends assumption, which
is supported by the absence of significant pre-1990 trends in wages or housing
values for station-adjacent tracts relative to control tracts.

The paper's preferred estimates show that flows between tract pairs where both
tracts contain a station rose by 11--16\,\%.  Our replication of the gravity
specification obtains a both-station effect of 10.0 log points (SE 3.8), stable
across eight progressively tighter geographic samples and confirmed by a Poisson
pseudo-maximum-likelihood (PPML) specification; adding the proximity ring within
250\,m raises the combined estimated effect to approximately 17.6\,\%.

\subsection{Structural elasticities}

To translate the reduced-form flow estimates into welfare-relevant parameters,
the paper estimates four structural elasticities using a Bartik shift-share IV
strategy.  The housing supply parameter $\psi$ is the most consequential.  It is
estimated by an IV regression of the change in housing prices
$(\Delta\ln P_{it})$ on the change in housing density $(\Delta\ln N_{it})$,
instrumenting density with a shift-share demand shock based on the 1990 industry
mix interacted with national industry-level employment growth.  The preferred
estimate is $\psi = 1.602$ (SE 0.453).  Note that $\psi$ is an \emph{inverse}
elasticity --- the slope of price on quantity --- so the implied housing supply
elasticity is $1/\psi = 0.624$: housing supply in LA is inelastic.

The first stage of the preferred specification has a Cragg--Donald $F$-statistic
of 11.5 and a Kleibergen--Paap $F$-statistic of 12.3; alternative specifications
in the paper's Table~4 range from 5.5 to 12.4.  The paper's bootstrap 95\,\%
confidence interval for $\psi$ is $[1.01, 3.67]$.  An internal consistency test
--- whether $\psi$ estimated from housing consumption and from residential land
use are consistent --- returns a $p$-value of 0.616, supporting the structural
interpretation.

\subsection{General-equilibrium welfare model}

The welfare calculation uses the hat-algebra approach of
\citet{dekle2008global}, extended to a multi-location setting with endogenous
commuting.  Workers choose a home and work location by drawing Fréchet-distributed
idiosyncratic preference shocks, giving rise to a gravity equation for commuting
flows:
%
\begin{equation}
  \pi_{ij} \;\propto\;
  \left(\frac{B_i\,W_j\,Q_i^{\,\zeta-1}}{\delta_{ij}}\right)^{\!\varepsilon},
  \label{eq:gravity}
\end{equation}
%
where $W_j$ is the wage at workplace tract $j$, $Q_i$ is the housing price at
residence tract $i$, $B_i$ captures residential amenities, $\delta_{ij}$ is the
bilateral commuting cost, $1-\zeta$ is the housing expenditure share, and
$\varepsilon$ is the Fréchet shape parameter governing the dispersion of
idiosyncratic location preferences.  Importantly, $\varepsilon$ is
\emph{estimated} in the paper (at $\approx 2.18$) using the same shift-share IV
strategy, rather than calibrated from the literature; the labour supply
first stage is in fact the weaker of the paper's IV designs (KP $F$ of 8.0
in the preferred specification).  The housing market clears at each tract with
inverse supply elasticity $\psi$.  The model is calibrated to the observed 2000
LA equilibrium and asked: by how much does welfare change when the observed
transit improvements are removed?  The answer --- solved by fixed-point
iteration converging in 6--11 iterations --- is the welfare gain attributable to
Metro Rail.  The baseline result is \textbf{\$93.55 million per year in
aggregate} (closed-city, 2016 dollars), matching the \$94 million reported in
the paper.\footnote{The dollar figure is the model's percentage welfare gain
(0.044\,\%) scaled by mean annual earnings (\$31{,}563) and the 6.73 million
workers in the five-county area.  It corresponds to roughly \$14 per worker per
year.}  This is the figure we take as our benchmark throughout.

% ============================================================
\section{Modifications and Methods}
\label{sec:methods}

\subsection{Motivation: heterogeneous housing supply}

\citet{severen2021} estimates a single pooled $\psi$ for the five-county LA metro
area.  The implicit assumption is that housing supply elasticity does not vary
systematically with station proximity or with the geographic characteristics of
the sub-market in which a station is located.  This assumption is strong:
coastal tracts face stricter development constraints than inland tracts;
already-dense central-city tracts have less room for upward densification than
low-density suburban tracts; and LA County faces very different zoning conditions
than Riverside or San Bernardino counties.

If $\psi$ is systematically \emph{higher} near stations --- supply more
inelastic, e.g.\ because stations were placed in already-dense,
supply-constrained areas --- the paper's aggregate welfare estimate understates
the extent to which transit benefits are capitalised into land values rather
than shared broadly.  Understanding the distribution of $\psi$ across LA's
sub-markets therefore matters for both efficiency and distributional
assessments of the Metro Rail programme.

\subsection{Attempted sub-regional estimation}
\label{subsec:subreg}

We attempt to estimate $\psi$ separately for three types of geographic
sub-regions, using the same Bartik IV specification as the original paper:

\begin{itemize}[leftmargin=1.5em, itemsep=2pt]
  \item \textbf{Coastal vs.\ inland.}  Tracts within 5\,km of the Pacific
    coastline versus all other tracts (296 coastal tracts,
    $\approx$2,257 inland).  Coastal distances are computed from a
    1:50\,000 Natural Earth coastline using California Albers projection.
  \item \textbf{Density quartiles.}  Tracts are ranked by 1990 homeowners
    per km$^2$ and split at the median (binary) and into quartiles.
    The prior is that already-dense tracts face less elastic supply
    (higher $\psi$) than low-density tracts.
  \item \textbf{County groups.}  A \emph{core} group (Los Angeles and Orange
    counties: dense, heavily zoned) versus a \emph{fringe} group (Riverside,
    San Bernardino, and Ventura counties: lower density and more
    developable land).  Both groups are large and industrially diverse,
    so the Bartik instrument was expected to retain adequate within-group
    variation.
\end{itemize}

Each specification uses the same IV as the original paper (the leave-one-out
Bartik employment shock $z_{it}$), restricts to the estimation sample
(\texttt{tc\_hval}~$= 0$), and weights by 1990 homeowners.

\subsection{Identification failure: instrument weakness}
\label{subsec:weakness}

Every sub-regional split encounters the same problem: the Bartik instrument
loses power when the sample is partitioned.  Table~\ref{tab:subreg} reports the
Kleibergen--Paap Wald $F$-statistic for each split.  The paper's pooled first
stage (KP $F = 12.3$ in the preferred specification) is moderately strong, and
our inland specification --- which retains $\approx$93\,\% of tracts and so
nearly proxies the full sample --- achieves KP $F = 9.89$.  But every genuine
partition falls to $F < 5$, and the coastal tracts, Density Q1, and the fringe
counties individually fall to $F < 1$.

The root cause is that the Bartik instrument derives its power from
cross-industry variation in employment growth interacted with the pre-period
industry mix.  When the sample is restricted to a geographic sub-area, the
effective number of tracts drops substantially, reducing the cross-sectional
variation available to identify the supply curve.  The coastal split is the
extreme case: with only 188 tracts, the instrument is essentially uninformative
(KP $F = 0.18$).

This exercise also calibrates how much identifying power the pooled design has
to spare: not much.  With KP $F$ statistics between roughly 7 and 12 across the
paper's own specifications and a wide bootstrap confidence interval for $\psi$
($[1.01, 3.67]$), the pooled estimate is adequately but not overwhelmingly
identified --- and there is clearly no surplus power for the sub-regional
heterogeneity analysis we attempted.

\subsection{GE sensitivity analysis}
\label{subsec:sensitivity}

Since sub-regional estimation is infeasible, we vary $\psi$ across a grid
spanning $0.5$ to $4.0$ and re-solve the full Severen GE model at each grid
point.  Because $\psi$ is an inverse elasticity, this grid spans highly
\emph{elastic} supply at the low end ($\psi = 0.5$ implies a supply elasticity
of $1/\psi = 2.0$, Sun Belt-type) to very \emph{inelastic} supply at the high
end ($\psi = 4.0$ implies $1/\psi = 0.25$, San Francisco-type).  All other
parameters are held at the paper's baseline values.  At each $\psi$ we verify
that:
%
\begin{itemize}[leftmargin=1.5em, itemsep=2pt]
  \item the contraction condition for the fixed-point iteration holds (largest
    eigenvalue of the transition matrix $< 1$);
  \item closed-city and open-city equilibria converge (typically in 6--11
    iterations to a tolerance of $10^{-5}$).
\end{itemize}
%
The eigenvalue ranges from 0.778 ($\psi = 0.5$) to 0.951 ($\psi = 4.0$), so the
iteration converges stably throughout.  (Stability of the iteration is distinct
from global uniqueness of the equilibrium; the model's separate parameter
condition for uniqueness is also satisfied at every grid point.)

\subsection{``Amplified-connectivity'' counterfactuals}
\label{sec:faster_transit}

As a complementary exercise, we ask how welfare would change if Metro Rail
generated a larger commuting benefit than actually observed (e.g.\ through
faster or more reliable service).

The commuting effect enters the model through the pair-specific bilateral
utility shifter $D_{ij}$.  In the hat-algebra counterfactual framework, the
change in $D_{ij}$ is governed by the estimated commuting parameters
$\lambda^{D}_{00}$ and $\lambda^{D}_{02}$, recovered from a
difference-in-differences regression of log bilateral commuting flows on binary
connectivity indicators $T^{00}_{ij}$ and $T^{02}_{ij}$---equal to one when both
home tract $i$ and workplace tract $j$ lie within 500\,m of a station of the
respective Metro line---controlling for OD-pair, residence-tract-by-year, and
workplace-tract-by-year fixed effects.

For multiplier $s \geq 1$, the counterfactual adjusts $D_{ij}$ as:
%
\begin{equation}
  \hat{D}_{ij}
  \;=\;
  \exp\!\Bigl(
    \lambda^{D}_{00}\,(s-1)\,T^{00}_{ij}
    \;+\;
    \lambda^{D}_{02}\,(s-1)\,T^{02}_{ij}
  \Bigr),
  \label{eq:Dhat}
\end{equation}
%
where $\lambda^{D}_{00} = 0.149$ and $\lambda^{D}_{02} = 0.128$.  Since both
parameters are positive, $\hat{D}_{ij} > 1$ for Metro-connected pairs at
$s > 1$, increasing their pair-specific commuting utility.  At $s = 1$ the
expression equals one (no change); at $s = 2$, each connected pair receives
$(s-1) = 1$ additional unit of the estimated commuting utility gain---doubling
the improvement attributed to Metro.  For OD pairs that are not
Metro-connected ($T^{00}_{ij} = T^{02}_{ij} = 0$ --- over 99.9\,\% of pairs,
carrying $\approx$99.2\,\% of commuters), $\hat{D}_{ij} = 1$ regardless of $s$.

It is worth noting that every Metro-connected pair receives the same
proportional utility shift, independent of actual trip length or actual travel
time, because the estimated $\lambda^D$ parameters capture a flat average
treatment effect.  The GE model is re-solved in full at each $s$, so downstream
price, wage, and population adjustments are endogenous.  This counterfactual
does not literally test trains running at a physically higher speed; it asks
what equilibrium would result if Metro connectivity raised commuting utility by
$s$ times its estimated magnitude.  We accordingly refer to $s$ as an
amplification multiplier, retaining ``speed'' only as informal shorthand.

Starting from the 2000 calibrated baseline, we add $(s-1)$ additional units of
the estimated benefit to all Metro-connected origin--destination pairs.  At
$s=2$ the commuting disutility reduction is twice its estimated magnitude.  The
GE model is re-solved at each multiplier and the incremental welfare gain on top
of the \$93.55 million baseline is recorded.

Re-solving the GE model across multipliers $s \in \{1.25, 1.5, 2, 3\}$ yields
welfare gains of \$25--\$223 million per year, increasing approximately linearly
in $s$.  This pattern is consistent with network coverage---only 94 of 2{,}552
tracts (4\,\%) lie near a station---rather than the magnitude of the per-pair
commuting benefit being the binding constraint on Metro Rail's welfare impact.
We note, however, that this comparison is indirect: we amplify the benefit on
existing corridors but do not simulate an expansion of the network itself, which
would be the direct test.

% ============================================================
\section{Results}
\label{sec:results}

\subsection{Aggregate welfare is insensitive to $\psi$}
\label{subsec:res_sensitivity}

Table~\ref{tab:psi_sensitivity} reports the closed-city welfare gain from Metro
Rail across the $\psi$ grid.  The headline finding is that aggregate welfare is
essentially invariant to $\psi$.  Moving from $\psi = 0.75$ to $\psi = 4.0$ ---
a more than $5\times$ increase --- reduces the welfare estimate by only \$0.33
million, from \$93.71 million to \$93.38 million.  This represents a 0.35\,\%
variation around the paper's baseline of \$93.55 million.  Over the full grid
(0.5 to 4.0), the variation is \$0.41 million.

The relationship is monotone: welfare is slightly \emph{higher} under more
elastic supply ($\psi = 0.5$, supply elasticity $2.0$: \$93.79 million) than
under more inelastic supply ($\psi = 4.0$, supply elasticity $0.25$: \$93.38
million).  The intuition is that with elastic supply, housing quantities near
stations can expand, letting more households take up the improved commutes;
with inelastic supply, more of the local adjustment is absorbed by price
increases, which in the closed-city model are transfers to incumbent landowners
rather than aggregate gains.  Either way the aggregate effect is tiny.  The
open-city population gain varies somewhat more (0.102\,\% at $\psi = 0.5$ down
to 0.081\,\% at $\psi = 4.0$): more elastic supply accommodates more
in-migration in response to the same transit improvement.

\begin{table}[H]
\centering
\caption{Welfare sensitivity to the inverse housing supply elasticity ($\psi$)}
\label{tab:psi_sensitivity}
\smallskip
\begin{tabular}{L{3.5cm} C{3.5cm} C{3cm} C{3cm}}
\toprule
\textbf{$\psi$} & \textbf{Closed-city welfare (\$M/year)} & \textbf{Gain (\% of income)} & \textbf{Open-city pop.\ gain} \\
\midrule
\shaderow 0.50            & \$93.79  & 0.0442\%  & 0.102\% \\
         0.75            & \$93.71  & 0.0441\%  & 0.097\% \\
\shaderow 1.00            & \$93.65  & 0.0441\%  & 0.093\% \\
         1.25            & \$93.60  & 0.0441\%  & 0.091\% \\
\shaderow 1.50            & \$93.56  & 0.0440\%  & 0.089\% \\
\textbf{1.602 (baseline)} & \textbf{\$93.55}  & \textbf{0.0440\%}  & \textbf{0.088\%} \\
\shaderow 1.75            & \$93.53  & 0.0440\%  & 0.087\% \\
         2.00            & \$93.50  & 0.0440\%  & 0.086\% \\
\shaderow 2.50            & \$93.46  & 0.0440\%  & 0.084\% \\
         3.00            & \$93.43  & 0.0440\%  & 0.083\% \\
\shaderow 4.00            & \$93.38  & 0.0440\%  & 0.081\% \\
\bottomrule
\end{tabular}

\medskip
\tablenote{GE model re-solved at each $\psi$ value; all other parameters held at
\citet{severen2021} baseline.  Baseline row in bold.  Welfare is aggregate
annual welfare for the five-county area in millions of 2016 dollars.  ``Gain
(\% of income)'' is the model's closed-city percentage welfare benefit
(compensating variation as a share of income).  Lower $\psi$ corresponds to more
elastic housing supply.  Source: \texttt{run\_welfare\_psi\_sensitivity.R}.}
\end{table}

\subsection{Sub-regional $\psi$ estimates: instrument weakness}
\label{subsec:res_subreg}

Table~\ref{tab:subreg} reports the sub-regional $\psi$ point estimates and
first-stage $F$-statistics.  The pattern is unambiguous: every geographic
partition that reduces the sample substantially pushes the Kleibergen--Paap
$F$-statistic below the conventional threshold of 10, and several produce $F <
1$.  The coastal split, with only 188 tracts, is completely uninformative ($F =
0.18$); the $\psi$ point estimate of 1.49 has a standard error of 5.72 and is
statistically indistinguishable from any value.  Even the county-group splits,
designed to preserve within-group industry diversity, produce $F$-statistics of
3.24 (core) and 1.93 (fringe).

The most nearly reliable sub-regional estimate is the inland specification
($\hat{\psi} = 0.49$, SE 0.17, KP $F = 9.89$), which proxies the full sample
because 93\,\% of tracts are inland.  Notably, this estimate is substantially
\emph{below} the paper's pooled $\psi = 1.602$, and the density-quartile
specifications where the instrument is marginally adequate (Q2, Q4) give
estimates in the range 0.54--0.59, also below 1.602.  Recalling that $\psi$ is
an inverse elasticity, lower values mean \emph{more elastic} housing supply: if
these marginally identified estimates were taken at face value, housing supply
in LA would be more elastic than the paper's pooled estimate implies.  Given
Section~\ref{subsec:res_sensitivity}, however, the aggregate welfare estimate
would be essentially unchanged in that case.

\begin{table}[H]
\centering
\caption{Sub-regional inverse housing supply elasticity ($\psi$) estimates}
\label{tab:subreg}
\smallskip
\begin{tabular}{L{6cm} C{1.6cm} C{1.6cm} C{1.6cm} C{1.8cm}}
\toprule
\textbf{Sub-region} & \textbf{$N$ tracts} & \textbf{$\hat{\psi}$} & \textbf{SE} & \textbf{KP $F$} \\
\midrule
\shaderow Full sample (pooled, Severen) & 2,175 & 1.602 & 0.453 & 12.3 \\
Core: LA + Orange Co. & $\approx$1,850 & 0.170 & 0.097 & 3.24 \\
\shaderow Fringe: Riverside + SB + Ventura & $\approx$590 & $-$0.650 & 0.489 & 1.93 \\
Coastal ($<$5\,km Pacific) & 188 & 1.488 & 5.716 & 0.18 \\
\shaderow Inland ($\geq$5\,km Pacific) & $\approx$2,254 & 0.486 & 0.169 & 9.89 \\
High density ($\geq$ median, 1990) & $\approx$1,200 & 0.435 & 0.174 & 1.83 \\
\shaderow Low density ($<$ median, 1990) & $\approx$1,200 & 0.493 & 0.400 & 7.46 \\
Density Q1 (least dense) & $\approx$600 & 0.088 & 1.080 & 0.007 \\
\shaderow Density Q2 & $\approx$600 & 0.590 & 0.311 & 6.48 \\
Density Q3 & $\approx$600 & 0.367 & 0.240 & 2.58 \\
\shaderow Density Q4 (most dense) & $\approx$600 & 0.540 & 0.255 & 6.47 \\
\bottomrule
\end{tabular}

\medskip
\tablenote{Each sub-regional row estimates $\psi$ by IV (\texttt{ivreg2}, Bartik
instrument, weighted by 1990 homeowners). KP $F$ = Kleibergen--Paap Wald
$F$-statistic.  Conventional threshold for adequate identification: $F \geq 10$.
Negative estimates for fringe counties reflect near-zero first stages and should
not be given structural interpretation.  Full-sample row reproduces the
preferred specification of \citet{severen2021}, Table~4 (column~6).
Source: \texttt{tracts\_elasticities\_coastal.do},
\texttt{tracts\_elasticities\_density.do},
\texttt{tracts\_elasticities\_countygroup.do}.}
\end{table}

\subsection{``Amplified-connectivity'' counterfactuals}
\label{subsec:res_faster}

Table~\ref{tab:faster} reports the welfare gains from amplifying the Metro Rail
commuting benefit, expressed as the incremental gain on top of the \$93.55
million baseline from current transit.  The relationship is roughly linear:
doubling the reduction in disutility ($s = 2\times$) adds \$105 million per
year, approximately four times the gain from a 25\,\% amplification (\$25
million).

\begin{table}[H]
\centering
\caption{Welfare gains from amplified Metro Rail connectivity}
\label{tab:faster}
\smallskip
\begin{tabular}{L{3.8cm} C{3.2cm} C{3.2cm} C{2.8cm}}
\toprule
\textbf{Multiplier $s$} & \textbf{Incremental gain (\$M/year)} & \textbf{Total welfare (\$M/year)} & \textbf{Open-city pop.\ gain} \\
\midrule
\shaderow 1.0$\times$ (current) & \$0   & \$93.55   & 0.000\% \\
         1.25$\times$           & \$25.17  & \$118.72  & 0.024\% \\
\shaderow 1.50$\times$          & \$51.09  & \$144.64  & 0.048\% \\
         2.00$\times$           & \$105.26 & \$198.81  & 0.099\% \\
\shaderow 3.00$\times$          & \$223.47 & \$317.02  & 0.211\% \\
         20.0$\times$           & \$5{,}921.70 & \$6{,}015.25 & --- \\
\bottomrule
\end{tabular}

\medskip
\tablenote{Incremental gain = additional aggregate welfare on top of current
transit (\$93.55 million/year).  The multiplier $s$ applies an additional
disutility reduction of $(s-1)\times$ the estimated effect on Metro corridors.
Open-city population gains are relative to the current-transit equilibrium.
The GE model converges at all levels of $s$.
Source: \texttt{simulate\_faster\_transit.R}; $20\times$ from
\texttt{map\_faster\_transit\_20x.R}.}
\end{table}

The $20\times$ counterfactual is extreme but instructive.  Aggregate welfare
reaches \$5{,}922 million per year, roughly 63 times the current transit benefit
and about 2.8\,\% of total annual labour earnings in the metro area.  Even at
this scale the response remains close to linear in $s$: within any plausible
range of transit improvements, the model's general-equilibrium feedbacks do not
generate non-linear amplification of the commuting benefit.

\subsection{The mechanism: why $\psi$ does not move aggregate welfare}
\label{subsec:mechanism}

The insensitivity of aggregate welfare to $\psi$ follows from two features of
the exercise: the small scale of the intervention, and the nature of the
closed-city welfare measure.

First, scale.  LA Metro Rail induces roughly a 10--18\,\% increase in commuting
flows between station pairs, but those pairs are a small fraction of all
$2{,}552 \times 2{,}552$ tract-pair combinations --- only 4\,\% of tracts are
near a station.  The aggregate demand shock to housing near stations is
therefore small relative to the total LA housing stock.

Second, what $\psi$ actually governs.  Monocentric city simulations (using
actual LA tract geography and an Alonso--Muth--Mills demand structure) make the
role of $\psi$ transparent.  For a local housing demand shock $D_n$ at tract
$n$, with demand elasticity $\eta$:
%
\begin{equation}
  \Delta\ln P_n \;=\; \frac{\psi\,D_n}{1 + \eta\,\psi}, \qquad
  \Delta\ln N_n \;=\; \frac{D_n}{1 + \eta\,\psi}.
  \label{eq:pq_split}
\end{equation}
%
A higher $\psi$ (more inelastic supply) pushes more of the local adjustment into
prices and less into quantities.  This split determines the \emph{incidence} of
the transit gains --- how much is capitalised into land values and captured by
incumbent landowners versus accommodated by new housing and shared with movers
--- but it contributes little to the \emph{aggregate} welfare number.  In the
closed-city measure, the price component is predominantly a transfer between
residents and incumbent landowners and largely nets out; the aggregate gain is
dominated by the direct commuting-cost reduction on connected pairs, which does
not involve $\psi$ at all.  Only when the demand shock is very large (as in the
$20\times$ counterfactual) do the GE price adjustments become large enough for
$\psi$ to have a visible aggregate effect.

This points to an important caveat on our own exercise.  Showing that aggregate
welfare is insensitive to $\psi$ does \emph{not} show that $\psi$ is irrelevant:
the distributional question that motivated this study --- who captures the gains
from transit --- is governed precisely by the price-vs-quantity split in
Equation~\eqref{eq:pq_split}.  Answering it would require decomposing the
model's gains by tract group and by owner/renter status, which the aggregate
welfare scalar cannot reveal.  We leave that decomposition to future work.

% ============================================================
\section{Conclusion}
\label{sec:conclusion}

This paper replicates the \citet{severen2021} analysis of LA Metro Rail welfare
and extends it by asking whether the city-wide assumption on housing supply
elasticity affects the results.  Our answer is that it does not affect the
aggregate welfare estimate, and we show why.

The attempt to estimate $\psi$ sub-regionally reveals a data limitation: the
Bartik instrument, moderately strong at the pooled level (KP $F = 12.3$ in the
preferred specification, with a wide bootstrap confidence interval for $\psi$),
loses power entirely when the sample is partitioned by geography or density.
This limits what can be learned about heterogeneity in housing supply conditions
across LA's sub-markets from this instrument and this dataset.

The GE sensitivity analysis, however, shows that this limitation is largely
inconsequential for the paper's headline welfare estimate.  Across a $5\times$
range of $\psi$, aggregate welfare moves by \$0.33 million --- 0.35\,\% of the
\$93.55 million baseline.  The mechanism is transparent: Metro Rail is a small
intervention in a large, car-dominated city, and the price adjustments that
$\psi$ governs are largely transfers within the closed city rather than
aggregate gains.  We emphasise that this is a statement about the aggregate:
$\psi$ continues to govern the incidence of the gains --- the split between
capitalisation into land values and broadly shared housing-quantity adjustment
--- which the aggregate welfare scalar cannot detect.

The amplified-connectivity counterfactuals add a policy implication.  Even
doubling the reduction in commuting disutility would add only \$105 million per
year, and the response is linear in the multiplier up to extreme values (at an
imaginary $20\times$ amplification the gain reaches \$5{,}922 million).  The
pattern is consistent with network coverage, rather than the per-pair magnitude
of the commuting benefit, being the binding constraint on Metro Rail's welfare
impact --- though a direct test would require simulating network expansion.

Taken together, these findings reinforce a central message of
\citet{severen2021}. The LA Metro Rail's welfare impact is modest because the city's
commuting patterns are fundamentally shaped by car infrastructure, and a limited
rail network serving 4\,\% of census tracts cannot substantially reshape spatial
equilibrium.  Housing supply elasticity, however it varies across the metro
area, is not the lever that would change this conclusion.

% ============================================================
\newpage
\bibliographystyle{plainnat}
\begin{thebibliography}{99}

\bibitem[Bartik(1991)]{bartik1991}
Bartik, T.~J. (1991).
\newblock \emph{Who Benefits from State and Local Economic Development
  Policies?}
\newblock W.~E. Upjohn Institute for Employment Research, Kalamazoo, MI.

\bibitem[Dekle et~al.(2008)Dekle, Eaton, and Kortum]{dekle2008global}
Dekle, R., Eaton, J., and Kortum, S. (2008).
\newblock Global rebalancing with gravity: Measuring the burden of adjustment.
\newblock \emph{IMF Staff Papers}, 55\penalty0 (3):\penalty0 511--540.

\bibitem[Diamond and McQuade(2019)]{diamond2019}
Diamond, R. and McQuade, T. (2019).
\newblock Who wants affordable housing in their backyard? An equilibrium
  analysis of low-income property development.
\newblock \emph{Journal of Political Economy}, 127\penalty0 (3):\penalty0
  1063--1117.

\bibitem[Goldsmith-Pinkham et~al.(2020)Goldsmith-Pinkham, Sorkin, and
  Swift]{goldsmith2020}
Goldsmith-Pinkham, P., Sorkin, I., and Swift, H. (2020).
\newblock Bartik instruments: What, when, why, and how.
\newblock \emph{American Economic Review}, 110\penalty0 (8):\penalty0
  2586--2624.

\bibitem[Monte et~al.(2018)Monte, Redding, and Rossi-Hansberg]{monte2018}
Monte, F., Redding, S.~J., and Rossi-Hansberg, E. (2018).
\newblock Commuting, migration, and local employment elasticities.
\newblock \emph{American Economic Review}, 108\penalty0 (12):\penalty0
  3855--3890.

\bibitem[Severen(2021)]{severen2021}
Severen, C. (2021).
\newblock Commuting, labor, and housing market effects of mass transportation:
  Welfare and identification.
\newblock \emph{Review of Economics and Statistics}, forthcoming.
\newblock \url{https://doi.org/10.1162/rest_a_01037}.

\end{thebibliography}

\end{document}
